\section*{Chapitre \ref{ch:m1:reading-market-data} --- Lecture des données de marché}

\begin{solution}{exo:m1:reading-market-data:1}
L'ordre 1001 disparaît (300 exécutés sur 300)~; l'ordre 1002 passe de 200 à 100.
\omterm{def:m1:reading-market-data:levels}{Niveau 2}~: cours acheteur 99.99 pour 100 en 1 ordre, 99.98 pour 100~; les cours vendeurs ne changent pas. \omterm{def:m1:reading-market-data:levels}{Niveau
1}~: 99.99 $\times$ 100 au cours acheteur, 100.01 $\times$ 500 au cours vendeur, dernier cours 99.99 $\times$ 300. La
transaction s'imprime à 99.99, le prix de l'ordre au repos, que le message d'exécution
ne transporte pas~; l'ordre au repos était un achat, donc l'agresseur était un
vendeur.
\end{solution}

\begin{solution}{exo:m1:reading-market-data:2}
\texttt{0x000F41DC} $= 999\,900$~: 99.9900~\$. Lu en little-endian, les mêmes octets
valent $3\,695\,251\,200$~: 369\,525.12~\$. Les erreurs d'ordre des octets sont généralement aussi
évidentes~; une erreur sur le nombre de décimales implicites (99.99 lu comme 9\,999) ne
l'est pas, dans un flux qui mêle des instruments.
\end{solution}

\begin{solution}{exo:m1:reading-market-data:3}
Chaque message coûte sa longueur plus 2 octets d'encadrement~: $1\,102 \times 38 + 311
\times 33 + 114 \times 25 + 473 \times 21 + 20 \times 46 = 41\,876 + 10\,263 + 2\,850 +
9\,933 + 920 = 65\,842$ octets.
\end{solution}

\begin{solution}{exo:m1:reading-market-data:4}
$2^{48}$ nanosecondes représentent 78.2 heures~: plus d'une journée, avec de la marge pour des sessions
qui chevauchent minuit. Quatre octets ne contiennent que 4.3 secondes de nanosecondes. Face à
huit octets, six en économisent deux sur chaque message, soit quelques pourcents d'un flux de
plusieurs dizaines de millions de messages par jour, au prix d'un champ non aligné que
chaque décodeur doit reconstituer à la main.
\end{solution}

\begin{solution}{exo:m1:reading-market-data:5}
Non. L'écart est de 40 microsecondes et chaque horloge peut être décalée de jusqu'à 100,
dans un sens ou dans l'autre~: l'ordre réel est indéterminé. Pour dire lequel a bougé
le premier, il faut une seule horloge pour les deux~: des \omterm{def:m1:reading-market-data:ts}{horodatages de réception} issus d'un même point
de capture, avec des latences connues et stables vers chaque plateforme, ou des plateformes synchronisées
un ordre de grandeur plus finement que l'écart mesuré.
\end{solution}

\begin{solution}{exo:m1:reading-market-data:6}
Arbitrage entre les deux lignes~: il prend 5\,001 sur la ligne qui le livre en premier, 5\,002
sur A, 5\,003 sur B, 5\,004 sur l'une ou l'autre, et le carnet est complet. Si les deux lignes
sautent 5\,003, le message est perdu~: demander une retransmission, ou attendre
le prochain instantané et reconstruire, en mettant les messages suivants en mémoire tampon dans l'intervalle. D'ici
là, le carnet des instruments concernés est marqué comme périmé~; la stratégie
doit cesser de coter à partir de lui, puisqu'elle ne connaît plus ses propres positions dans la file
ni même le meilleur prix, et ne peut réduire son risque qu'à l'aide d'une
source de prix plus lente et indépendante.
\end{solution}

\begin{solution}{exo:m1:reading-market-data:7}
Environ 1~\% à 100 microsecondes et 11~\% à 1 milliseconde. Lorsque chaque second
message est perdu, les écarts entre messages consécutifs doublent~; une même gigue dépasse donc
moins de voisins, et les pourcentages tombent à environ la moitié
pour de petites gigues. Le réordonnancement dépend de la gigue \emph{relative au
temps entre messages}~: il est pire dans les rafales, justement quand l'ordre compte
le plus.
\end{solution}

\begin{solution}{exo:m1:reading-market-data:8}
(i) Mauvais ticks~: les erreurs de décimales et les transactions annulées, que les barres du fournisseur
contiennent encore, produisent des «~chutes~» de 50~\% ou plus qui se résorbent à la barre suivante~;
personne n'aurait pu acheter à ce niveau. (ii) Conditions~: des comptes rendus tardifs et hors séquence
ou des impressions issues du niveau de prix d'une autre journée, entrent dans le plus bas. (iii)
\omterm{def:m1:reading-market-data:refdata}{Symbologie}~: des barres indexées par symbole réunissent deux sociétés lorsqu'un symbole est réutilisé,
et des splits non ajustés ressemblent à une chute de 50~\% à l'ouverture. On peut ajouter
le biais de survie dans la liste des symboles, ainsi que le fait qu'un plus bas de barre ne dit rien
sur la taille disponible à ce prix.
\end{solution}

\begin{solution}{pb:m1:reading-market-data:1}
\textbf{1.} 22\,650 actions à 98.40.
\textbf{2.} 10.01.
\textbf{3.} $98.40/100.05 - 1 = -165$ points de base.
\textbf{4.} Le calcul explose~: deux rendements consécutifs de $-90\%$ et $+900\%$ dominent
toute estimation pendant des mois.
\textbf{5.} Une erreur de décimale ou de saisie~: un cours dix fois trop petit entre
deux impressions normales à une seconde d'intervalle. Un filtre sur la distance à une médiane glissante
ou à la cotation en vigueur le détecte~; il faut le confirmer par la cotation à
09:30:30 (aucun cours acheteur proche de 10 n'existait) et à l'annulation ultérieure
de la bourse.
\textbf{6.} Son prix est normal. Seule sa \omterm{def:m1:reading-market-data:condition}{condition de transaction}, ou un message d'annulation distinct
renvoyant à son \omterm{def:m1:reading-market-data:seq}{numéro de séquence}, la retire~: étape 1 de la
\cref{met:m1:reading-market-data:clean}.
\textbf{7.} Non. Il s'agit du volume du jour, qui figure dans un total journalier, mais il
n'a pas été exécuté pendant la vie de l'ordre et son prix reflète 09:24~; un indice de référence
pour 09:30 à 09:32 doit utiliser l'heure d'exécution, et non l'heure du
compte rendu.
\textbf{8.} Dedans~: c'est une transaction réelle dans l'intervalle, à un prix de marché.
(Savoir si les \omterm{def:m1:us-equity-market-structure:lots}{lots irréguliers} mettent à jour le dernier cours consolidé relève d'une règle de publication~;
pour une moyenne pondérée par les volumes, ils comptent.)
\textbf{9.} Dehors, pour ce client~: l'enchère a clos avant que l'ordre n'existe
et il ne pouvait pas y participer. Il serait dedans pour un ordre saisi avant
l'ouverture.
\textbf{10.} Trois exclues (impressions 6, 9 et 11)~; 16\,950 actions~; 100.045.
\textbf{11.} 4\,950 actions à 100.106.
\textbf{12.} Plus bas 100.02 (l'enchère) ou 100.04 en négociation continue~; plus haut
100.15.
\textbf{13.} Par rapport à 100.045~: $+4.5$ points de base (il a payé plus cher). Par rapport à
100.106~: $-1.6$ (il a payé moins cher).
\textbf{14.} La seconde. Un indice de référence doit contenir les transactions auxquelles l'ordre aurait pu
participer~; l'enchère représente 71~\% du volume du premier chiffre
et était terminée avant l'arrivée de l'ordre.
\textbf{15.} Sous forme de message d'annulation ou de rupture plus tard dans la journée, ou dans le
fichier de corrections du lendemain, renvoyant à la transaction d'origine par son \omterm{def:m1:reading-market-data:seq}{numéro de séquence} ou
numéro de correspondance. Le chargeur doit l'appliquer à la transaction stockée, conserver les deux
enregistrements et recalculer tout ce qui en dérive.
\textbf{16.} Parce qu'il existe toujours un ensemble d'exclusions qui rend une
exécution donnée bonne. Une règle choisie après avoir vu le résultat n'est pas un
indice de référence.
\textbf{17.} La liste exacte des conditions exclues, le traitement des
corrections et des comptes rendus tardifs, l'horodatage utilisé pour l'affectation aux barres, et
la question de savoir si les \omterm{def:m1:us-equity-market-structure:lots}{lots irréguliers} et les transactions hors bourse sont inclus.
\textbf{18.} Capture~: rien n'est écarté, tout est conservé. \omterm{def:m1:reading-market-data:norm}{Normalisation}~: les conditions
sont transformées en indicateurs explicites, les corrections sont liées aux originaux. Construction du carnet~:
non concernée pour les transactions, mais la cotation en vigueur qu'elle produit est ce qui
détecte l'impression 6. Tables~: la règle d'inclusion propre à chaque usage et l'indicateur de
valeur aberrante.
\textbf{19.} \textbf{Trois impressions exclues~; 100.045 (100.106 pour la négociation continue
uniquement).}
\textbf{20.} Un message signalant qu'un intervenant a déclaré une transaction, avec un prix, une
taille, une heure et un ensemble de qualifications qui déterminent ce qu'elle signifie.
\end{solution}

\begin{solution}{iq:m1:reading-market-data:1}
\omterm{def:m1:reading-market-data:levels}{Niveau 1}~: meilleurs cours acheteur et vendeur, avec leurs tailles, et dernière transaction. \omterm{def:m1:reading-market-data:levels}{Niveau 2}~: taille totale
à chaque prix, par niveau. \omterm{def:m1:reading-market-data:levels}{Niveau 3}~: chaque ordre et chaque événement le concernant. Chacun se
calcule à partir du suivant~; seul le \omterm{def:m1:reading-market-data:levels}{niveau 3} donne la position dans la file, les durées de vie
des ordres et l'identité de ce qui a été exécuté.
\iqlookfor{la relation à sens unique et ce que seul le \omterm{def:m1:reading-market-data:levels}{niveau 3} fournit.}
\end{solution}

\begin{solution}{iq:m1:reading-market-data:2}
Une table de hachage qui associe la référence de chaque ordre à celui-ci (sens, prix, taille restante)
et, par instrument et par sens, une structure triée qui associe chaque prix à la taille agrégée
(et, pour la position dans la file, les ordres à ce prix dans l'ordre chronologique). Ajout~:
insertion dans les deux. Exécution ou annulation~: recherche par référence, réduction des deux, retrait
si le niveau est vide. Suppression~: idem avec la taille complète. Les meilleurs prix sont aux extrémités de
la structure triée. Les carnets de production remplacent les arbres par des tableaux indexés
par tick autour du meilleur cours, et pré-allouent la mémoire.
\iqlookfor{recherche par référence, car les exécutions ne portent aucun prix.}
\end{solution}

\begin{solution}{iq:m1:reading-market-data:3}
Vérifier d'abord l'autre ligne. Si les deux l'ont manqué~: marquer les carnets concernés
comme obsolètes, prévenir les stratégies, mettre en mémoire tampon ce qui suit, demander une retransmission ou
prendre l'instantané suivant, reconstruire, rejouer la mémoire tampon, lever l'indicateur. Consigner le
trou avec leurs heures. Négocier sur un carnet obsolète est interdit~; retirer ses cotations est
le comportement par défaut.
\iqlookfor{le drapeau d'obsolescence parvenant à la stratégie.}
\end{solution}

\begin{solution}{iq:m1:reading-market-data:4}
L'\omterm{def:m1:reading-market-data:ts}{horodatage de réception}, au point où la stratégie aurait vu les
données, auquel s'ajoute la latence propre de la stratégie pour agir. Les \omterm{def:m1:reading-market-data:ts}{horodatages de la bourse} donnent
au backtest une information antérieure à son arrivée~: il négocie sur des cotations déjà disparues, et les
signaux entre plateformes paraissent rentables parce que l'horloge d'une
plateforme avance sur celle de l'autre. L'\omterm{def:m1:reading-market-data:ts}{horodatage de la bourse} reste pertinent pour
ordonner les événements au sein d'une même plateforme.
\iqlookfor{biais d'anticipation par l'horloge.}
\end{solution}

\begin{solution}{iq:m1:reading-market-data:5}
Identifiants~: ce qui sert de clé au fichier, et s'ils restent stables dans le temps. Heure~:
quelle horloge, quel fuseau, heure d'exécution ou de déclaration. Conditions~: lesquelles sont
présentes, lesquelles ont été retirées. Corrections~: appliquées ou non. Unités~: devise, échelle des
prix, ajustées ou brutes pour les \omterm{def:m1:shares-corporate-actions-indices:split}{divisions d'actions}. Doublons et trous~: comptages par jour, à comparer à
un total indépendant. Valeurs aberrantes par rapport aux cotations. Couverture~: quelles
plateformes, en bourse et hors bourse. Puis dix minutes à examiner une journée à l'œil.
\iqlookfor{conditions, corrections, identifiants, horloge.}
\end{solution}

\begin{solution}{iq:m1:reading-market-data:6}
Conserver la capture brute, sans modification, avec les \omterm{def:m1:reading-market-data:ts}{horodatages de réception} matériels,
compressée et dotée de sommes de contrôle, par plateforme, ligne et jour~: c'est la seule
vérité de référence. Versionner le décodeur et les \omterm{def:m1:reading-market-data:refdata}{données de référence}~; consigner, pour chaque table dérivée,
les fichiers de capture, la version du code et les paramètres qui l'ont produite.
Les données dérivées (carnets, barres, tables) sont des caches que l'on peut supprimer et
reconstruire. Les \omterm{def:m1:reading-market-data:refdata}{données de référence} sont conservées telles qu'elles étaient à chaque date, sans jamais être écrasées. Tester
la reproductibilité en reconstruisant chaque semaine une ancienne journée tirée au hasard et en comparant
les empreintes de hachage.
\iqlookfor{capture brute comme vérité de référence~; \omterm{def:m1:reading-market-data:refdata}{données de référence} historisées~;
provenance de chaque table.}
\end{solution}
